\section{Dual Spaces and Dual Maps}\label{sec:dual-spaces}
Coordinates record a vector after a basis has been chosen.  There is another
useful way to probe a vector: apply every scalar-valued linear measurement to
it.  This viewpoint is intrinsic, and it is the starting point for
bilinear forms and alternating forms.

\begin{definition}
A \emph{linear functional} on an \(F\)-vector space \(V\) is a linear map
\(\varphi:V\to F\).  The \emph{algebraic dual space} of \(V\) is
\[
V^\vee=\operatorname{Hom}_F(V,F).
\]
\end{definition}

The notation \(V^\vee\) is deliberately different from the notation
\(T^*\) used for the inner-product adjoint in Chapter~7. The construction
below is algebraic and requires no inner product.

\begin{proposition}[Dual basis]\label{prop:dual-basis}
Let \(\beta=(v_1,\ldots,v_n)\) be an ordered basis of a finite-dimensional
vector space \(V\).  There are unique linear functionals
\[
v_i^\vee:V\longrightarrow F
\]
such that
\[
v_i^\vee(v_j)=\delta_{ij}.
\]
Moreover,
\[
\beta^\vee=(v_1^\vee,\ldots,v_n^\vee)
\]
is a basis of \(V^\vee\).  In particular,
\[
\dim V^\vee=\dim V.
\]
\end{proposition}
\begin{proof}
For
\[
v=\sum_{j=1}^na_jv_j,
\]
define \(v_i^\vee(v)=a_i\).  The uniqueness of basis coordinates proves
that this is well-defined and linear, and gives the displayed values.
If
\[
\sum_{i=1}^nc_iv_i^\vee=0,
\]
evaluation at \(v_j\) gives \(c_j=0\), so the functionals are independent.
For \(\varphi\in V^\vee\), linearity gives
\[
\varphi=\sum_{i=1}^n\varphi(v_i)v_i^\vee.
\]
Thus they span and form a basis.
\end{proof}

Consequently \(V\) and \(V^\vee\) are isomorphic when \(V\) is
finite-dimensional, but this proof used the chosen basis.  In general there
is no preferred isomorphism
\[
V\cong V^\vee.
\]
An inner product can supply one, but that is additional structure.

\begin{definition}
For a linear map
\[
T:V\longrightarrow W,
\]
its \emph{dual map}, also called its algebraic transpose, is
\[
T^\vee:W^\vee\longrightarrow V^\vee,
\qquad
T^\vee(\varphi)=\varphi\circ T.
\]
\end{definition}

The direction reverses because a measurement on \(W\) can be pulled back
along \(T\) to a measurement on \(V\).

\begin{proposition}[Functorial identities and transpose matrix]
For linear maps \(T:V\to W\) and \(S:W\to U\),
\[
(\operatorname{id}_V)^\vee=\operatorname{id}_{V^\vee},
\qquad
(S\circ T)^\vee=T^\vee\circ S^\vee.
\]
If \(\beta\) and \(\gamma\) are ordered bases of \(V\) and \(W\), then
\[
[T^\vee]^{\beta^\vee}_{\gamma^\vee}
=\left([T]^\gamma_\beta\right)^t.
\]
\end{proposition}
\begin{proof}
For \(\varphi\in V^\vee\),
\[
(\operatorname{id}_V)^\vee(\varphi)
=\varphi\circ\operatorname{id}_V=\varphi.
\]
Likewise, for \(\psi\in U^\vee\),
\[
(S\circ T)^\vee(\psi)
=\psi\circ S\circ T
=T^\vee(S^\vee(\psi)).
\]
Write
\[
T(v_j)=\sum_{i=1}^ma_{ij}w_i,
\qquad
[T]^\gamma_\beta=(a_{ij}).
\]
For the \(i\)-th dual basis vector of \(W^\vee\),
\[
T^\vee(w_i^\vee)(v_j)=w_i^\vee(T(v_j))=a_{ij}.
\]
Hence
\[
T^\vee(w_i^\vee)=\sum_{j=1}^na_{ij}v_j^\vee.
\]
Its matrix therefore has \((j,i)\)-entry \(a_{ij}\), which is the transpose
of the matrix of \(T\).
\end{proof}

\begin{definition}
The \emph{evaluation map} is
\[
\iota_V:V\longrightarrow V^{\vee\vee},
\qquad
\iota_V(v)(\varphi)=\varphi(v).
\]
\end{definition}

\begin{proposition}[Double dual]\label{prop:double-dual}
The evaluation map \(\iota_V\) is linear and injective.  If \(V\) is
finite-dimensional, it is an isomorphism.
\end{proposition}
\begin{proof}
For \(a,b\in F\), \(v,w\in V\), and \(\varphi\in V^\vee\),
\[
\iota_V(av+bw)(\varphi)
=\varphi(av+bw)
=a\iota_V(v)(\varphi)+b\iota_V(w)(\varphi),
\]
so \(\iota_V\) is linear.  If \(v\ne0\), extend \(v\) to a basis and take
the corresponding first coordinate functional.  It does not vanish on
\(v\), so \(\iota_V(v)\ne0\).  Thus \(\iota_V\) is injective.  In finite
dimension, Proposition~\ref{prop:dual-basis} gives
\[
\dim V^{\vee\vee}=\dim V,
\]
so an injective map between these finite-dimensional spaces is surjective.
\end{proof}

Thus finite-dimensional vector spaces have a canonical isomorphism
\[
V\cong V^{\vee\vee},
\]
even though an isomorphism with the first dual is usually noncanonical.
This canonical character is expressed by the naturality identity
\begin{equation}\label{eq:double-dual-naturality}
\iota_W\circ T=T^{\vee\vee}\circ\iota_V
\qquad(T:V\to W).
\end{equation}
Indeed, for \(v\in V\) and \(\varphi\in W^\vee\),
\[
\begin{aligned}
\bigl(T^{\vee\vee}\iota_V(v)\bigr)(\varphi)
&=\iota_V(v)(\varphi\circ T)\\
&=(\varphi\circ T)(v)\\
&=\varphi(Tv)\\
&=\iota_W(Tv)(\varphi).
\end{aligned}
\]
Thus evaluation commutes with every linear map and requires no basis choice.

\begin{definition}
For a subspace \(W\leq V\), its \emph{annihilator} is
\[
W^\circ=\{\varphi\in V^\vee:\varphi(w)=0\text{ for every }w\in W\}.
\]
\end{definition}

\begin{proposition}
If \(V\) is finite-dimensional, then
\[
\dim W^\circ=\dim V-\dim W.
\]
\end{proposition}
\begin{proof}
Choose a basis \(w_1,\ldots,w_r\) of \(W\) and extend it to a basis
\[
(w_1,\ldots,w_r,u_1,\ldots,u_s)
\]
of \(V\).  Its dual basis shows that \(W^\circ\) has basis
\[
(u_1^\vee,\ldots,u_s^\vee).
\]
Since \(r+s=\dim V\), the result follows.
\end{proof}

\begin{theorem}[Dual maps, kernels, and images]\label{thm:dual-map-kernel-image}
Let \(T:V\to W\) be a linear map between finite-dimensional vector spaces.
Then
\[
\ker(T^\vee)=(\operatorname{im}T)^\circ,
\qquad
\operatorname{im}(T^\vee)=(\ker T)^\circ.
\]
Consequently,
\[
\operatorname{rank}(T^\vee)=\operatorname{rank}(T).
\]
\end{theorem}
\begin{proof}
For \(\varphi\in W^\vee\), every equivalence is direct:
\[
\begin{aligned}
\varphi\in\ker(T^\vee)
&\Longleftrightarrow T^\vee\varphi=0\\
&\Longleftrightarrow \varphi\circ T=0\\
&\Longleftrightarrow \varphi\text{ vanishes on }\operatorname{im}T\\
&\Longleftrightarrow \varphi\in(\operatorname{im}T)^\circ.
\end{aligned}
\]
Hence, using the annihilator dimension formula and rank--nullity,
\[
\begin{aligned}
\operatorname{rank}(T^\vee)
&=\dim W^\vee-\dim\ker(T^\vee)\\
&=\dim W-\bigl(\dim W-\dim\operatorname{im}T\bigr)\\
&=\operatorname{rank}T.
\end{aligned}
\]
If \(\varphi=T^\vee\psi=\psi\circ T\) and \(v\in\ker T\), then
\(\varphi(v)=\psi(Tv)=0\).  Thus
\(\operatorname{im}(T^\vee)\subseteq(\ker T)^\circ\).  The two spaces have
the same dimension, because
\[
\dim\operatorname{im}(T^\vee)=\operatorname{rank}T
=\dim V-\dim\ker T=\dim(\ker T)^\circ.
\]
The containment is therefore equality.
\end{proof}

In dual bases,
\([T^\vee]=[T]^t\).  The theorem therefore yields
\[
\operatorname{rank}(A^t)=\operatorname{rank}(A),
\]
giving the conceptual dual-space explanation of the row-rank equals
column-rank theorem proved earlier by elementary matrix methods.

\begin{proposition}[Quotient dual]\label{prop:quotient-dual}
If \(W\leq V\) and \(\pi:V\to V/W\) is the quotient map, then
\[
\pi^\vee:(V/W)^\vee\longrightarrow V^\vee
\]
is injective with image \(W^\circ\).  Hence
\[
(V/W)^\vee\cong W^\circ.
\]
\end{proposition}
\begin{proof}
Because \(\pi\) is surjective, \(\psi\circ\pi=0\) implies \(\psi=0\), so
\(\pi^\vee\) is injective.  Theorem~\ref{thm:dual-map-kernel-image} gives
\[
\operatorname{im}(\pi^\vee)=(\ker\pi)^\circ=W^\circ.
\]
\end{proof}

\begin{proposition}[Annihilator identities]\label{prop:annihilator-identities}
For subspaces \(U,W\leq V\), with \(V\) finite-dimensional,
\[
U\subseteq W\Longrightarrow W^\circ\subseteq U^\circ,
\qquad
(U+W)^\circ=U^\circ\cap W^\circ,
\]
\[
(U\cap W)^\circ=U^\circ+W^\circ,
\qquad
W^{\circ\circ}=W
\]
under the evaluation identification \(V\cong V^{\vee\vee}\).
\end{proposition}
\begin{proof}
A functional vanishing on \(W\) certainly vanishes on \(U\subseteq W\),
which proves reversal of inclusions.  A functional vanishes on \(U+W\)
exactly when it vanishes on both summands, proving the second identity.

Every element of \(U^\circ+W^\circ\) vanishes on \(U\cap W\), so
\[
U^\circ+W^\circ\subseteq(U\cap W)^\circ.
\]
The dimension formula for sums and intersections gives
\[
\begin{aligned}
\dim(U^\circ+W^\circ)
&=\dim U^\circ+\dim W^\circ-\dim(U^\circ\cap W^\circ)\\
&=(\dim V-\dim U)+(\dim V-\dim W)-\dim(U+W)^\circ\\
&=\dim V-\dim(U\cap W),
\end{aligned}
\]
which is the dimension of \((U\cap W)^\circ\); hence equality holds.
Finally, evaluation sends every \(w\in W\) to a functional on \(V^\vee\)
that vanishes on \(W^\circ\), so \(W\subseteq W^{\circ\circ}\).  Both spaces
have dimension \(\dim W\), and therefore they are equal.
\end{proof}

\section{Bilinear and Quadratic Forms}\label{sec:bilinear-quadratic-forms}
Duality turns a vector into a linear measurement.  A bilinear form instead
assigns a scalar to an ordered pair of vectors, linearly in each argument.
This is the algebraic framework behind dot products, but it does not include
positivity or conjugation.

\begin{definition}
A \emph{bilinear form} on an \(F\)-vector space \(V\) is a map
\[
B:V\times V\longrightarrow F
\]
that is linear in each variable.  It is \emph{symmetric} if
\[
B(x,y)=B(y,x),
\]
\emph{skew-symmetric} if
\[
B(x,y)=-B(y,x),
\]
and \emph{alternating} if
\[
B(v,v)=0
\qquad(v\in V).
\]
\end{definition}

Every alternating bilinear form is skew-symmetric, by expanding
\[
0=B(x+y,x+y).
\]
Conversely, a skew-symmetric form is alternating only when the
characteristic is not \(2\), in the sense of
Definition~\ref{def:ring-characteristic}. This distinction is essential in
characteristic \(2\).
In every ordered basis \(\beta\),
\[
B\text{ symmetric}\Longleftrightarrow[B]_\beta^t=[B]_\beta,
\qquad
B\text{ skew-symmetric}\Longleftrightarrow[B]_\beta^t=-[B]_\beta.
\]
Alternating always implies the second condition.  If
\(\operatorname{char}F\ne2\), the second condition forces
\(2B(v,v)=0\) and hence \(B(v,v)=0\).  In characteristic \(2\), however,
\(-A=A\), so the transpose equation merely says that \(A\) is symmetric; an
alternating form must additionally have zero diagonal.

\begin{proposition}[Matrix and congruence]\label{prop:bilinear-matrix}
For an ordered basis \(\beta=(v_1,\ldots,v_n)\), put
\[
[B]_\beta=(B(v_i,v_j)).
\]
Then
\[
B(x,y)=[x]_\beta^t[B]_\beta[y]_\beta.
\]
If
\[
P=[I_V]^\beta_{\beta'},
\]
so that
\[
[x]_\beta=P[x]_{\beta'},
\]
then
\[
[B]_{\beta'}=P^t[B]_\beta P.
\]
\end{proposition}
\begin{proof}
Write
\[
x=\sum_i a_iv_i,
\qquad
y=\sum_j b_jv_j.
\]
By bilinearity,
\[
B(x,y)=\sum_{i,j}a_iB(v_i,v_j)b_j,
\]
which is the stated matrix product.  The columns of \(P\) are the
\(\beta\)-coordinates of the vectors of \(\beta'\).  Substitution of
\[
[x]_\beta=P[x]_{\beta'},
\qquad
[y]_\beta=P[y]_{\beta'}
\]
into the coordinate formula gives
\[
B(x,y)
=[x]_{\beta'}^tP^t[B]_\beta P[y]_{\beta'},
\]
and uniqueness of a matrix representing the form in \(\beta'\) proves the
claim.
\end{proof}

This is \emph{congruence}, not similarity.  For an operator, one input is
converted into output coordinates and change of basis gives
\[
[T]_{\beta'}=P^{-1}[T]_\beta P.
\]
A bilinear form has two vector inputs, so both coordinates are substituted
and the formula has \(P^t\) on the left instead.

\begin{definition}
The \emph{left radical} and \emph{right radical} of \(B\) are, respectively,
\[
\operatorname{rad}_L(B)
=\{v\in V:B(v,w)=0\text{ for every }w\in V\},
\]
\[
\operatorname{rad}_R(B)
=\{w\in V:B(v,w)=0\text{ for every }v\in V\}.
\]
For a general bilinear form these are conceptually distinct.  The associated
linear maps are
\[
\Phi_B:V\longrightarrow V^\vee,
\qquad
\Phi_B(v)=B(v,-).
\]
and
\[
\Psi_B:V\longrightarrow V^\vee,
\qquad
\Psi_B(w)=B(-,w).
\]
\end{definition}

\begin{definition}
The \emph{rank of a bilinear form} is
\[
\operatorname{rank}B:=\operatorname{rank}\Phi_B.
\]
\end{definition}

\begin{proposition}[Rank and radicals of a bilinear form]
For every ordered basis \(\beta\),
\[
\operatorname{rank}B=\operatorname{rank}[B]_\beta.
\]
This number is independent of the basis.  Moreover,
\[
\dim\operatorname{rad}_L(B)=\dim V-\operatorname{rank}B,
\qquad
\dim\operatorname{rad}_R(B)=\dim V-\operatorname{rank}B.
\]
\end{proposition}
\begin{proof}
In \(\beta\) and its dual basis, the matrix of \(\Phi_B\) is
\([B]_\beta^t\).  Equality of row and column rank therefore gives
\[
\operatorname{rank}\Phi_B
=\operatorname{rank}[B]_\beta^t
=\operatorname{rank}[B]_\beta.
\]
Alternatively, basis independence follows directly from
\([B]_{\beta'}=P^t[B]_\beta P\) and invariance of rank under multiplication
by invertible matrices.  Rank--nullity applied to \(\Phi_B\) proves the
left-radical formula.  The right-radical formula follows in the same way from
\(\Psi_B\), whose matrix is \([B]_\beta\).
\end{proof}

\begin{proposition}[Nondegeneracy criterion]
For a finite-dimensional vector space, the following are equivalent:
\begin{enumerate}
\item \(\operatorname{rad}_L(B)=\{0\}\);
\item \(\operatorname{rad}_R(B)=\{0\}\);
\item \([B]_\beta\) is invertible for one, equivalently every, basis
\(\beta\).
\end{enumerate}
When these conditions hold, \(B\) is called \emph{nondegenerate}.
\end{proposition}
\begin{proof}
By their definitions,
\[
\ker\Phi_B=\operatorname{rad}_L(B),
\qquad
\ker\Psi_B=\operatorname{rad}_R(B).
\]
In the basis \(\beta\) and its dual, the matrices of these maps are
\[
[\Phi_B]_{\beta}^{\beta^\vee}=[B]_\beta^t,
\qquad
[\Psi_B]_{\beta}^{\beta^\vee}=[B]_\beta.
\]
Indeed, the \(j\)-th coordinate of \(\Phi_B(v_i)\) is \(B(v_i,v_j)\),
whereas that of \(\Psi_B(v_i)\) is \(B(v_j,v_i)\).  Thus either map is
injective exactly when \([B]_\beta\) is invertible.  Since
\[
\dim V=\dim V^\vee,
\]
injectivity is equivalent to bijectivity for both maps.  Finally, the
congruence formula in Proposition~\ref{prop:bilinear-matrix} shows that
matrix invertibility is independent of the chosen basis.
\end{proof}

Equivalently, in finite dimension,
\[
\begin{aligned}
B\text{ nondegenerate}
&\Longleftrightarrow \operatorname{rank}B=\dim V\\
&\Longleftrightarrow [B]_\beta\text{ is invertible}
\Longleftrightarrow \det[B]_\beta\ne0.
\end{aligned}
\]
A chosen nondegenerate bilinear form therefore supplies the isomorphism
\[
\Phi_B:V\xrightarrow{\ \sim\ }V^\vee,\qquad v\longmapsto B(v,-).
\]
This is precisely the extra structure missing from the bare vector space:
there is no preferred isomorphism \(V\cong V^\vee\), but a nondegenerate form
provides one.

For a symmetric bilinear form, \(B(v,w)=B(w,v)\), so the two radicals
coincide.  In that case we write their common value as \(\operatorname{rad}(B)\).

\begin{definition}
For a symmetric bilinear form \(B\) and a subspace \(W\leq V\), define
\[
W^{\perp_B}=\{v\in V:B(v,w)=0\text{ for every }w\in W\}.
\]
\end{definition}

\begin{proposition}[Orthogonal complements for a bilinear form]
If \(B\) is nondegenerate on finite-dimensional \(V\), then
\[
\dim W^{\perp_B}=\dim V-\dim W.
\]
Moreover,
\[
V=W\oplus W^{\perp_B}
\quad\Longleftrightarrow\quad
B|_W\text{ is nondegenerate}.
\]
\end{proposition}
\begin{proof}
The isomorphism \(\Phi_B:V\to V^\vee\) carries \(W^{\perp_B}\) onto the
annihilator \(W^\circ\).  The annihilator dimension formula gives the first
claim.  Also,
\[
W\cap W^{\perp_B}=\operatorname{rad}(B|_W).
\]
Thus the sum is direct exactly when \(B|_W\) is nondegenerate; in that case
the dimension formula shows that the direct sum has dimension \(\dim V\) and
therefore equals \(V\).  The converse is immediate from the displayed
intersection.
\end{proof}

Unlike the positive-definite case, nondegeneracy of \(B\) on \(V\) alone does
not guarantee this direct sum.  For example, on \(F^2\) with
\[
[B]=\begin{pmatrix}0&1\\1&0\end{pmatrix},
\]
the form is nondegenerate, but \(e_1\) is isotropic and
\(Fe_1=(Fe_1)^{\perp_B}\).  Thus isotropic vectors can make
\(W\cap W^{\perp_B}\ne0\).

\begin{definition}
Assume \(\operatorname{char}F\ne2\), in the sense of
Definition~\ref{def:ring-characteristic}. A quadratic form associated with a
symmetric bilinear form \(B\) is
\[
q(v)=B(v,v).
\]
\end{definition}

The form can be recovered from its quadratic function by polarization:
\[
B(x,y)
=\frac12\bigl(q(x+y)-q(x)-q(y)\bigr).
\]
Characteristic \(2\) requires a separate theory: a quadratic form must not
there be identified blindly with a symmetric bilinear form.

In coordinates, if \(A=[B]_\beta=A^t\) and
\(x=(x_1,\ldots,x_n)^t\), then
\[
q(x)=x^tAx
=\sum_i a_{ii}x_i^2+2\sum_{i<j}a_{ij}x_ix_j.
\]
Conversely, the polarization formula recovers the unique symmetric bilinear
form from \(q\).  If \(x=Py\) and \(P^tAP=D\), then
\[
q(Py)=(Py)^tA(Py)=y^tP^tAPy=y^tDy
=d_1y_1^2+\cdots+d_ny_n^2.
\]
Thus congruence is exactly the coordinate change that diagonalizes a
quadratic form.

\begin{example}[Completing the square]
Over a field of characteristic not \(2\), consider
\[
q(x,y)=x^2+4xy+3y^2.
\]
The cross term disappears visibly:
\[
q(x,y)=(x+2y)^2-y^2.
\]
Set \(u=x+2y\) and \(v=y\), equivalently
\[
\begin{pmatrix}x\\y\end{pmatrix}
=\begin{pmatrix}1&-2\\0&1\end{pmatrix}
\begin{pmatrix}u\\v\end{pmatrix}.
\]
Then \(q=u^2-v^2\).  This is the coordinate mechanism behind diagonalization,
not a classification of all quadratic forms.
\end{example}

\begin{theorem}[Diagonalization of symmetric bilinear forms]
Suppose that \(\operatorname{char}F\ne2\).  Every symmetric bilinear form on
a finite-dimensional \(F\)-vector space has a basis in which its matrix is
diagonal.  Equivalently, its quadratic form has coordinates
\[
q(x_1,\ldots,x_n)=a_1x_1^2+\cdots+a_nx_n^2.
\]
\end{theorem}
\begin{proof}
Induct on \(\dim V\).  If \(B=0\), there is nothing to prove.  Otherwise
some \(v\) satisfies \(B(v,v)\ne0\): if every such value were zero, then
polarization would give \(B=0\).  Let
\[
W=\{w\in V:B(v,w)=0\}.
\]
For \(x\in V\),
\[
x=\frac{B(x,v)}{B(v,v)}v+
\left(x-\frac{B(x,v)}{B(v,v)}v\right),
\]
and the parenthesized vector lies in \(W\).  Thus
\[
V=Fv\oplus W.
\]
The restriction of \(B\) to \(W\) is symmetric, so induction supplies a
diagonalizing basis of \(W\); adjoining \(v\) gives the desired basis.
\end{proof}

Over \(\mathbb R\), inner products are precisely positive-definite symmetric
bilinear forms. Over \(\mathbb C\), the inner products of Chapter~7 are
Hermitian: they are linear in the first variable and conjugate-linear in the
second. Thus Chapter~7 is not merely the complex special case of this
bilinear discussion.
